% Part: turing-machines
% Chapter: machines-computations
% Section: halting-states

\documentclass[../../../include/open-logic-section]{subfiles}

\begin{document}

\olfileid{tur}{mac}{hal}
\olsection{停机状态}

\begin{explain}
尽管我们将机器定义为仅在没有指令可执行时才停机，但在图灵机的常见表示中，都会包含一个专门的\emph{停机状态}~$h$，使得 $h \in Q$。

停机状态的思想很简单：当机器完成操作（准备好接收输入，或已完成输出写入）时，它就会进入状态~$h$ 并停机。有些机器有两个停机状态，一个用于接受输入，一个用于拒绝输入。
\end{explain}

\begin{ex}\emph{停机状态}。
为了阐明这一概念，让我们从修改偶数机开始。与其让机器在输入为偶数时停机于状态~$q_0$，我们可以增加一条指令，让机器进入一个停机状态。
\[
\begin{tikzpicture}[->,>=stealth',shorten >=1pt,auto,node distance=2.8cm,
                    semithick]
  \tikzstyle{every state}=[fill=none,draw=black,text=black]

  \node[initial, state]         (A)                     {$q_0$};
  \node[state]         (B) [right of=A] {$q_1$};
  \node[state]         (C) [below of=A] {$h$};

  \path (A) edge [bend left] node {\TMtrans{\TMstroke}{\TMstroke}{R}} (B)
  	    edge node {\TMtrans{\TMblank}{\TMblank}{N}} (C)
        (B) edge [loop above] node {\TMtrans{\TMblank}{\TMblank}{R}} (B)
            edge [bend left] node {\TMtrans{\TMstroke}{\TMstroke}{R}} (A);
\end{tikzpicture}
\]

让我们进一步扩展这个例子。当机器判定输入为奇数时，它永不停机。我们可以修改机器使其包含一个\emph{拒绝}状态，方法是将循环指令替换为一条转入拒绝状态~$r$ 的指令。
\[
\begin{tikzpicture}[->,>=stealth',shorten >=1pt,auto,node distance=2.8cm,
                    semithick]
  \tikzstyle{every state}=[fill=none,draw=black,text=black]

  \node[initial,state]         (A)                     {$q_0$};
  \node[state]         (B) [right of=A] {$q_1$};
  \node[state]         (C) [below of=A] {$h$};
  \node[state]         (D) [below of=B] {$r$};

  \path (A) edge [bend left] node {\TMtrans{\TMstroke}{\TMstroke}{R}} (B)
  	    edge node {\TMtrans{\TMblank}{\TMblank}{N}} (C)
        (B) edge node {\TMtrans{\TMblank}{\TMblank}{N}} (D)
            edge [bend left] node {\TMtrans{\TMstroke}{\TMstroke}{R}} (A);
\end{tikzpicture}
\]
\end{ex}

\begin{explain}
在这种情况下，增加一个专用停机状态是有好处的，因为它明确了机器何时接受/拒绝某些输入。但需要注意的是，增加专用停机状态并不会增强计算能力。类似地，较不形式化的停机概念也有其自身的优势。本章迄今为止使用的停机定义，使得\emph{停机问题}的证明直观且易于展示。因此，我们继续沿用最初的定义。
\end{explain}

\end{document}